\section{Limitations}
\label{sec:limitations}
This study covers four consumer systems and a narrow Qwen3 model family. Hardware, operating systems, runtime versions, backends, and some quantization choices co-vary. The observed failures of local indicators are case evidence, not an estimated population-wide failure rate or a controlled portability benchmark. The diagnosed mechanisms are not exhaustive and do not causally explain the M4 workload contrast. The results establish observed regimes and within-platform mechanisms, not an architectural ranking or a general deployment law. Serving concurrency, production load, and broader model families are outside the measured scope.

Token equivalence is deliberately strict and conditional on the frozen reference; a different quality contract could admit other paths. HumanEval functional execution and pass@$k$ were not evaluated. Serialized B570 outputs still mismatch, and the CUDA evidence is bounded to the captured stack and case. No claim of general backend unsuitability follows.

Campaign-specific estimators and bootstrap implementations differ. Prompt intervals do not cover all system-state uncertainty, including M4 session/thermal effects. The M4 corrected synthetic-ceiling preregistration remains unexecuted. Decode timing excludes initial prefill and loading, so it is not whole-request user latency. Energy sensor domains differ; there is no global energy ranking or end-to-end 358H Phase-3 energy claim.

Windows raw per-run provenance is weaker than that of the other closed campaigns. The 358H draft-source reconstruction is \texttt{STRONGLY\_SUPPORTED}, not proven historical upstream identity. Replay and proposal are measured together, preventing separate causal timing attribution to either. No measured runtime intervention exists in this version; persistent state, path selection, and profitability control remain proposals. These limitations bound the hardware conclusion independently of the future software decision.
